% ===================================================================
%  TABELLA N-RO 1 — Le schola. — Le lyceo. — Le classe.
%
%  Ceci n'est PAS une transcription : c'est une traduction, faite en
%  2026, en interlingua d'aujourd'hui. D'ou trois differences avec
%  texto/io et texto/fr, et trois seulement :
%
%   * pas de \nl ni de \cc — la traduction ne suit aucune ligne
%     imprimee, et les lignes de ce fichier ne sont que du confort ;
%   * pas de \begin{VUpage} — elle n'a ni feuillet ni folio, et la
%     page de lecture ne lui en affiche pas ;
%   * le gras se pose sur le TERME seul, ponctuation dehors.
%
%  Tout le reste est identique, et doit l'etre : les cles « %%K » sont
%  celles de texto/io, macro pour macro pour l'apparat, et les renvois
%  \textsuperscript{(n)} visent les MEMES objets de la planche — ce
%  sont eux qui ouvrent les gros plans.
%
%  LA DEUXIEME DES LANGUES QUE LES IDISTES PARLENT DEJA. L'interlingua
%  de l'IALA parait en 1951, un quart de siecle apres ce livret ; elle
%  ne descend pas de l'ido mais du meme projet — donner a l'Europe une
%  langue commune tiree de son fonds latin. On traduit donc SUR L'IDO,
%  comme pour l'esperanto.
%
%  MAIS ELLE NE SE FABRIQUE PAS, ELLE SE RELEVE. L'esperanto compose
%  ses mots par affixes — « lernejo », « instruisto » ; l'interlingua
%  prend celui que les langues sources ont deja : « schola »,
%  « professor », « fenestra », « pupitro ». La ou l'esperanto dirait
%  « klasĉambro », l'interlingua dit « classe », parce que c'est le
%  mot que le francais, l'espagnol, l'italien, l'anglais et l'allemand
%  ont en commun. Le mot ne se derive pas ; il se constate.
%
%  NI ACCUSATIF NI ACCORD, contrairement a l'esperanto : l'adjectif de
%  l'interlingua est invariable — « le tabellas auxiliar », « duo
%  belle vasos ». C'est ce qui la rapproche le plus de l'ido, et c'est
%  pourquoi l'ordre des renvois se garde sans effort.
%
%  LES NOMS DE PERSONNE SONT NATURALISES, comme dans les quatorze
%  autres colonnes : Ioannes est Johannes, Iakobus est Jacobo, Karolus
%  est Carolo. Seuls le chien Medor et l'ane Aliboron gardent leur nom.
%  Les noms latins de la note (*) ne bougent pas : c'est d'eux que la
%  note parle.
% ===================================================================

%%K t01-apar-0 apar
\VUsaut{2.0mm}
\VUcentre{12.6pt}{120}{LIBRETTO EXPLICATIVE}
\VUsaut{3.2mm}
\VUcentre{8.4pt}{160}{DEL}
\VUsaut{2.6mm}
\VUtitre{15.6pt}{0}{Tabellas Auxiliar Delmas}
\VUsaut{3.4mm}
\VUornamento{9pt}
\VUsaut{5.0mm}
\VUcentre{13.2pt}{90}{PRIME SERIE}
\VUsaut{3.0mm}
\VUfilet{16mm}
\VUsaut{5.6mm}

%%K t01-apar-1 apar
\VUtitre{15.0pt}{60}{TABELLA N\textsuperscript{o} 1}
\VUsaut{1.4mm}
\VUtitre{11.4pt}{0}{Le schola, le lyceo, le classe.}
\VUsaut{2.2mm}
\VUfilet{20mm}
\VUsaut{6.4mm}

%%K t01-tit-1 sub
\VUpk{10.2pt}{40}{Description general.}
\VUsaut{3.0mm}

%%K t01-01-1 p
1. --- Iste tabella representa un \VUgras{classe} in un \VUgras{lyceo}. In iste
classe il ha octo \VUgras{tabulas} \textsuperscript{(18)} e octo \VUgras{bancos}
\textsuperscript{(32)}. Sur cata banco sede tres alumnos. Le \VUgras{professor}
\textsuperscript{(7)} sede sur un \VUgras{sedia} \textsuperscript{(8)} in su
\VUgras{cathedra} \textsuperscript{(6)}.

%%K t01-02-1 p
2. --- Al sinistra del cathedra se trova un \VUgras{armario}
\textsuperscript{(15)} \VUgras{vitrate}. Al dextra es le \VUgras{tabula nigre}
\textsuperscript{(1)} con le \VUgras{spongia} \textsuperscript{(2)} e le
\VUgras{creta} \textsuperscript{(3)}.

%%K t01-02-2 p
Super le cathedra pende \VUgras{tabellas mural} \textsuperscript{(9, 11, 12)} e
un \VUgras{carta} \textsuperscript{(10)}.

%%K t01-03-1 p
3. --- Non tote le alumnos es a lor \VUgras{loco} \textsuperscript{(17)}.
Johannes \textsuperscript{(4)} sta al tabula, ille tene le spongia e dele le
\VUgras{figuras geometric}.

%%K t01-03-2 p
Petro es presso le cathedra, ille collige le \VUgras{devores scripte}
\textsuperscript{(5)}, e los porta al professor.

%%K t01-04-1 p
4. --- Iste professor es le professor del \VUgras{linguas vive}. Ille insenia
al alumnos parlar, leger e scriber linguas estranier \VUgras{(le germano, le
anglese, le espaniol, le italiano, le russo} o \VUgras{le francese)}.

%%K t01-05-1 p
5. --- Le classe es multo vaste e multo clar. In le fundo il ha un large
\VUgras{fenestra} con duo \VUgras{fenestrettas} \textsuperscript{(78)} e le
\VUgras{porta vitrate} pro aerar e illuminar lo.

%%K t01-05-2 p
Iste classe non es frigide. In le medio il ha, pro calefacer lo, un multo bon
\VUgras{stufa} \textsuperscript{(46)} con su \VUgras{tubo}
\textsuperscript{(45)}.

%%K t01-05-3 p
Al pariete es fixate \VUgras{crocos pro vestimentos} \textsuperscript{(58)}, de
illos pende le cappellos e le mantellos del scholares.

%%K t01-tit-2 sub
\VUsaut{5.2mm}
\VUpk{10.2pt}{40}{Le corte de recreation}
\VUsaut{3.6mm}

%%K t01-06-1 p
6. --- Le \VUgras{horologio} \textsuperscript{(63)} indica novem horas e cinque
minutas.

%%K t01-06-2 p
Le \VUgras{recreation} \textsuperscript{(76)} justo finiva e le lection
recomencia. Le loco de recreation se trova in le \VUgras{corte}
\textsuperscript{(75)}. Nos vide alcun alumnos qui juga. Un \VUgras{sub-maestro}
\textsuperscript{(74)} los survelia.

%%K t01-07-1 p
7. --- In le corte nos vide ancora diverse personas, nominatemente le
\VUgras{inspector} \textsuperscript{(73)}, le director del schola
\VUgras{(director del lyceo)} \textsuperscript{(71)} e le \VUgras{censor}
\textsuperscript{(72)}.

%%K t01-07-2 p
Le inspector inspecta le scholas, le director del lyceo dirige le lyceo, le
censor survelia le studios.

%%K t01-07-3 p
Le quarte persona es le \VUgras{portero} \textsuperscript{(77)}. Ille porta sub
le bracio un registro pro inscriber le absentes.

%%K t01-08-1 p
8. --- In le fundo del corte il ha le \VUgras{apparatos gymnastic}
\textsuperscript{(70)} e le officina pro le \VUgras{labores manual}
\textsuperscript{(69)}.

%%K t01-08-2 p
Al dextra del tabella nos comencia a vider le \VUgras{classes} de
\VUgras{musica} e de \VUgras{designo}.

%%K t01-noto-1 noto
\VUnotes{143.2mm}{%
(*) Pro le \textit{nomines de baptismo} on consilia anque transscriber los
secundo le forma latin. Isto es le unic medio pro escappar a divergentias
innumerabile. Del resto como eliger inter \textit{Giovanni, Jean,}
\textit{Johann, Jan, Hans, John, Ivan,} etc.? Va nos crear un dictionario
special pro le nomines de baptismo? Prendamos del latino lor nominativo e
dicamos : \VUgras{Ioannes, Ioanna; Iulius, Iulia; Iakobus; Andreas; Lukas.}
(Gram. Complete).%
}

%%K t01-tit-3 sub
\VUpk{10.2pt}{40}{Description detaliate.}
\VUsaut{2.4mm}

%%K t01-tit-4 sub
\VUpk{10.2pt}{40}{Le bon alumnos.}
\VUsaut{3.0mm}

%%K t01-09-1 p
9. --- Le alumnos del prime banco al dextra del cathedra es \VUgras{Henrico}
(*), \VUgras{Stephano} e \VUgras{Carolo}.

%%K t01-09-2 p
Henrico es multo attentive e laboriose, ille ascolta le professor.

Sur su tabula ille ha un \VUgras{quaderno} \textsuperscript{(28)}, un
\VUgras{libro} \textsuperscript{(29)}, un \VUgras{lexico}
\textsuperscript{(30)} e un \VUgras{portaplumas} \textsuperscript{(31)}. Su
libro ha multe \VUgras{paginas}, cata capitulo contine plure
\VUgras{paragraphos} e cata \VUgras{linea} multe \VUgras{parolas}.

%%K t01-09-4 p
Su vicino Stephano \textsuperscript{(24)} es fervorose. Ille nota sur su
quaderno le lection pro le vice sequente. Ille ha un belle
\VUgras{scriptura}. Su libro es claudite. Nos vide solmente le \VUgras{copertura}
\textsuperscript{(27)}. Presso ille se trova su \VUgras{compasso}
\textsuperscript{(26)}.

%%K t01-09-5 p
Carolo \textsuperscript{(23)} tene con le mano su libro, ille lo aperi pro
releger su lection. Ille ha, como cata alumno, un \VUgras{incastro}
\textsuperscript{(25)} ante se. Ille es obediente.

%%K t01-10-1 p
10. --- Al tabula vicin es \VUgras{Ludovico, Antonio} e \VUgras{Paulo}. Ludovico
recita su \VUgras{lection} \textsuperscript{(22)} e responde al
\VUgras{questiones} del professor. Antonio \textsuperscript{(21)} es multo
inattentive, mesmo alcun vices le professor le puni. Paulo
\textsuperscript{(20)} es un bon camerada, affabile e complacente. Ille es
multo attentive.

%%K t01-tit-5 sub
\VUsaut{4.2mm}
\VUpk{10.2pt}{40}{Le mal alumnos.}
\VUsaut{3.2mm}

%%K t01-11-1 p
11. --- Detra Henrico se trova \VUgras{Georgio} \textsuperscript{(38)}. Ille non
es un mal puero, ma ille es \VUgras{garrule}, inattentive e petulante. Su
\VUgras{frivolitate} e \VUgras{disobedientia} causa \VUgras{disordine} durante
le lection e sovente ille merita un sever \VUgras{advertimento}. Su
\VUgras{quaderno de brossa} \textsuperscript{(35)} es immunde, ille lo
\VUgras{macula de incausto} con su \VUgras{penna} \textsuperscript{(34)}, pro
isto ille usa sempre su \VUgras{gumma pro deler} \textsuperscript{(36)}; su
\VUgras{portaquadernos} \textsuperscript{(33)} es lacerate, proque ille es multo
negligente. Proque porta ille su \VUgras{portacrayon} \textsuperscript{(37)} in
le classe del linguas vive?

%%K t01-11-2 p
Su vicino Edmundo \textsuperscript{(39)} non le ama. Ille es in ver un pauco
indolente, ma ille es taciturne e tranquille. Ille non garrula; ma, in loco de
ascoltar, ille prefere leger le \VUgras{narrationes} de su \VUgras{libro de
lectura}.

%%K t01-11-3 p
Le tertie alumno de iste banco es \VUgras{Ludovico} \textsuperscript{(40)}. Ille
es diligente. Ille scribe sur un blanc \VUgras{folio de papiro}. Ante se, ille
poneva su \VUgras{papiro absorbente} \textsuperscript{(41)}.

%%K t01-12-1 p
12. --- Le alumnos del secunde banco, al sinistra del professor, es multo juvene.
Le prime, le parve \VUgras{Emilio}, non sape ancora multo ben scriber, ille
porta su \VUgras{ardesia} \textsuperscript{(42)}, su \VUgras{stilo de ardesia} e
su \VUgras{spongia} \textsuperscript{(43)}.

%%K t01-12-2 p
Presso ille es \VUgras{Augusto} e \VUgras{Bernardo} \textsuperscript{(44)}. Iste
ultime labora ben, ma su \VUgras{vestimento} es multo negligite.

%%K t01-13-1 p
13. --- Detra illes se trova tres \VUgras{indolentes}. \VUgras{Alberto} non
apprende su lectiones. \VUgras{Jacobo} \textsuperscript{(47)} dormi in loco de
ascoltar. Su \VUgras{indolentia} le vale \VUgras{reproches} e mesmo
\VUgras{punitiones}. Finalmente \VUgras{Christiano} \textsuperscript{(48)} es
troppo frivole. Ille \VUgras{se comporta} mal e non se effortia in nulle modo
pro emendar se.

%%K t01-14-1 p
14. --- Su vicinos, al contrario, es de bon conducta. \VUgras{Ernesto}
\textsuperscript{(51)} ama le ordine e le regularitate, ille usa sempre su
\VUgras{regula} \textsuperscript{(50)} e su \VUgras{crayon}
\textsuperscript{(49)}. Ille es \VUgras{externo}.

%%K t01-14-2 p
\VUgras{Francisco} \textsuperscript{(52)} es \VUgras{interno}, ille es vestite
de un uniforme, mangia e dormi in le lyceo. Ille es un bon \VUgras{camerada}.

%%K t01-14-3 p
Mauricio talia su \VUgras{crayon} con su \VUgras{cultelletto}
\textsuperscript{(56)}, ille es multo curose.

%%K t01-14-4 p
Ille porta tote su libros, su \VUgras{grammatica} \textsuperscript{(55)}, su
\VUgras{carnet} \textsuperscript{(54)}, e mesmo un \VUgras{cossino pro scriber}
\textsuperscript{(53)}. Su \VUgras{sacco de schola} \textsuperscript{(57)} es
sur le solo detra ille.

%%K t01-14-5 p
Le duo tabulas in le fundo del classe es occupate per \VUgras{bon alumnos}
\textsuperscript{(19)}.

%%K t01-tit-6 sub
\VUsaut{4.6mm}
\VUpk{10.2pt}{40}{Designo e musica.}
\VUsaut{3.4mm}

%%K t01-15-1 p
15. --- In le \VUgras{sala de designo} il ha gratiose \VUgras{modellos}
appendite al muro e un \VUgras{lampa} \textsuperscript{(62)} con
\VUgras{umbrella}, suspendite al plafon. Le \VUgras{professor}
\textsuperscript{(60)} es multo patiente, ille corrige le \VUgras{designos}
\textsuperscript{(61)} del alumnos e les insenia designar un \VUgras{busto}
\textsuperscript{(59)} secundo le natura.

%%K t01-16-1 p
16. --- Le \VUgras{sala de musica} pare multo interessante. In illo on apprende
leger le \VUgras{musica}, solfear le \VUgras{notas del gamma}
\textsuperscript{(67)} scribite sur le \VUgras{pentagramma} (do, re, mi, fa,
sol, la, si\,=\,C. D. E. F. G. A. B.), cognoscer le \VUgras{claves} e cantar
gratiose \VUgras{melodias}. On pone le partituras sur le \VUgras{pupitro}
\textsuperscript{(65)}. Un del alumnos \VUgras{sona le piano}
\textsuperscript{(64)} e le \VUgras{professor de musica} \textsuperscript{(66)}
\VUgras{batte le mesura} \textsuperscript{(68)}.

%%K t01-17-1 p
17. --- Nos comencia a vider un angulo del \VUgras{officina} pro \VUgras{labores
manual} \textsuperscript{(69)}, ubi le pueros pote apprender planar le ligno e
limar le ferro. Isto non existe in tote le lyceos.

%%K t01-tit-7 sub
\VUsaut{5.0mm}
\VUpk{10.2pt}{40}{Le sala de scientias.}
\VUsaut{3.6mm}

%%K t01-18-1 p
18. --- Le professor de mathematica insenia al alumnos le \VUgras{geometria, le
arithmetica, le calculo} e le \VUgras{systema metric}. Nos vide sur le tabella
le principal figuras del geometria : \VUgras{circulo} \textsuperscript{(a)},
\VUgras{quadrato} \textsuperscript{(b)}, \VUgras{triangulo}
\textsuperscript{(c)}, \VUgras{linea} \textsuperscript{(d)}.

%%K t01-18-2 p
Nos vide anque un \VUgras{calculo}; illo non es ni un \VUgras{addition}, ni un
\VUgras{subtraction}, ni un \VUgras{multiplication} : illo es un
\VUgras{division} \textsuperscript{(e)}.

%%K t01-19-1 p
19. --- Sur le tabella del systema metric nos vide : un \VUgras{metro}
\textsuperscript{(a)}, un \VUgras{balancia} \textsuperscript{(b)} e su
\VUgras{pesos}, un \VUgras{litro} \textsuperscript{(e)}, un \VUgras{decalitro}
\textsuperscript{(f)}, un \VUgras{decilitro} e un \VUgras{metro cubic}
\textsuperscript{(d)}.

%%K t01-19-2 p
Le alumnos studia anque le \VUgras{scientias natural} \textsuperscript{(11)} ---
le zoologia \textsuperscript{(a)}, le botanica \textsuperscript{(b)}, le
geologia \textsuperscript{(c)} --- e le \VUgras{historia}
\textsuperscript{(12)}.

%%K t01-19-3 p
In le armario il ha le apparatos de \VUgras{physica} \textsuperscript{(15)} e de
\VUgras{chimia} \textsuperscript{(16)}.

%%K t01-19-4 p
Super le armario il ha : un \VUgras{sphera} \textsuperscript{(14)}, un
\VUgras{cono} e un \VUgras{globo terrestre} \textsuperscript{(13)}.

%%K t01-19-5 p
Super le tabellas pende un \VUgras{carta} que monstra le cinque partes del
mundo : \VUgras{Europa} \textsuperscript{(g)}, \VUgras{Asia}
\textsuperscript{(e)}, \VUgras{Africa} \textsuperscript{(f)}, \VUgras{America}
\textsuperscript{(ab)} e \VUgras{Oceania} \textsuperscript{(h)}, e le duo grande
oceanos, le \VUgras{Pacific} \textsuperscript{(c)} e le \VUgras{Atlantic}
\textsuperscript{(d)}.
\VUsaut{6.0mm}
\VUfilet{22mm}
